\section*{Chapter \ref{ch:m3:forward-curves-storage-and-convenience-yield} --- Forward Curves, Storage and Convenience Yield}

\begin{solution}{exo:m3:forward-curves-storage-and-convenience-yield:1}
$80\,e^{(0.04 + 0.03 - 0.05)} = 80\,e^{0.02} = \$81.62$.
\end{solution}

\begin{solution}{exo:m3:forward-curves-storage-and-convenience-yield:2}
\omterm{def:m3:forward-curves-storage-and-convenience-yield:carry}{Full carry}: $80\,(e^{0.07/12} - 1) = \$0.47$ a barrel.
\end{solution}

\begin{solution}{exo:m3:forward-curves-storage-and-convenience-yield:3}
A contango steeper than \omterm{def:m3:forward-curves-storage-and-convenience-yield:carry}{full carry} is closed by the \omterm{def:m3:forward-curves-storage-and-convenience-yield:carry}{cash-and-carry trade}, which anyone with storage can
do. Arbitraging backwardation requires selling the commodity now and buying it back later, which only
holders of inventory can do, and they hold it because they need it: the \omterm{def:m3:forward-curves-storage-and-convenience-yield:storage}{convenience yield} has no upper
bound.
\end{solution}

\begin{solution}{exo:m3:forward-curves-storage-and-convenience-yield:4}
$0.05 - 12\ln(71.20/70.00) = -15.4\%$ a year: the curve is in contango by more than financing, so the
market is paying for storage; \omterm{def:m3:forward-curves-storage-and-convenience-yield:storage}{convenience yield} is near zero and inventories are ample.
\end{solution}

\begin{solution}{exo:m3:forward-curves-storage-and-convenience-yield:5}
$5.55 - 3.08 - 3.18 = -0.71\%$ a year.
\end{solution}

\begin{solution}{exo:m3:forward-curves-storage-and-convenience-yield:6}
One day: $0.004 \times 50 \times 1 = \$0.20$ a barrel. Ten days: $0.004 \times 5 \times 5.5 = \$0.11$.
Spreading the roll lowers the cost but lengthens the window others can trade ahead of.
\end{solution}

\begin{solution}{exo:m3:forward-curves-storage-and-convenience-yield:7}
40.9\%, 38.3\%, 35.3\% and 33.6\%; the ratio of the fourth to the first is 0.82. With $e^{-\kappa \times
0.25} = 0.82$, $\kappa = 0.79$ a year (a half-life of about 0.88 years).
\end{solution}

\begin{solution}{exo:m3:forward-curves-storage-and-convenience-yield:8}
The fund holds futures, not oil: it earns the spot change plus the \omterm{def:m3:forward-curves-storage-and-convenience-yield:roll}{roll yield} plus interest. With the
curve in contango more than half the time, the roll cost 0.71\% a year from 1985 to 2024, and more in
some decades; fees and the index's own roll costs come on top.
\end{solution}

\begin{solution}{pb:m3:forward-curves-storage-and-convenience-yield:1}
\textbf{1.} 10 million barrels. \textbf{2.} \$0.04. \textbf{3.} $0.04 \times 3 = \$0.12$ a barrel.
\textbf{4.} USD~6 million. \textbf{5.} USD~72 million. \textbf{6.} It sells the second contract and buys
the third, the same spread the index will trade. \textbf{7.} During the window, selling the spread back
to the index: it earns the spread's rise, up to \$0.20 a barrel on positions taken before the first day.
\textbf{8.} Capital, margin, the risk that the spread moves for other reasons, and competition from other
front-runners. \textbf{9.} Early positioning moves the spread before the window: the index still pays,
but the gain is shared among more traders. \textbf{10.} The traders buy and sell at market prices on
public information; they take the other side of a known trade, which is liquidity provision in advance.
\textbf{11.} \$0.20 in one day, \$0.11 over ten. \textbf{12.} Transparency and replicability: products
tracking the index must reproduce it exactly. \textbf{13.} Deferred contracts sit where the curve is
flatter, so the \omterm{def:m3:forward-curves-storage-and-convenience-yield:roll}{roll yield} in contango is less negative. \textbf{14.} By comparing spread returns inside
and outside the roll period, 2000 to March 2010: statistically significant price impact, front-running
Sharpe ratios as high as 4.39, and about 3.6\% of annual return lost by index investors. \textbf{15.} A
lower return than the spot suggests, from both \omterm{def:m3:forward-curves-storage-and-convenience-yield:roll}{roll yield} and roll cost. \textbf{16.} Both: it makes the
index replicable and cheap to run, but charges the investor the impact. \textbf{17.} In contango the new
contract is dearer than the old: the position rolls into more expensive contracts every month.
\textbf{18.} The curve's shape (\omterm{def:m3:forward-curves-storage-and-convenience-yield:roll}{roll yield}), the index's roll rule and its impact, fees, and the \omterm{def:m3:forward-curves-storage-and-convenience-yield:roll}{collateral
return}. \textbf{19.} \emph{Named result:} \$0.12 a barrel per roll, USD~6 million a month and USD~72
million a year in the illustrative model. \textbf{20.} Whoever provides the storage or the inventory the
curve pays for: holders of the commodity in backwardation, storers in contango.
\end{solution}

\begin{solution}{iq:m3:forward-curves-storage-and-convenience-yield:1}
The implicit benefit of holding the physical commodity: the option to use it now. It is high when
inventories are low and a stock-out is costly, which shows as backwardation.
\iqlookfor{inventory dependence and the curve's shape.}
\end{solution}

\begin{solution}{iq:m3:forward-curves-storage-and-convenience-yield:2}
\omterm{def:m3:forward-curves-storage-and-convenience-yield:roll}{Spot return} (the nearby price's change), \omterm{def:m3:forward-curves-storage-and-convenience-yield:roll}{roll yield} (convergence of held contracts, positive in
backwardation) and \omterm{def:m3:forward-curves-storage-and-convenience-yield:roll}{collateral return} (interest on the margin cash).
\iqlookfor{three parts, and the sign of \omterm{def:m3:forward-curves-storage-and-convenience-yield:roll}{roll yield}.}
\end{solution}

\begin{solution}{iq:m3:forward-curves-storage-and-convenience-yield:3}
Inventories and supply respond to shocks over time, so shocks to near delivery are larger than to far:
the \omterm{def:m3:forward-curves-storage-and-convenience-yield:samuelson}{Samuelson effect}. Model the spot as mean-reverting (Ornstein--Uhlenbeck) plus a slow factor: each
future's volatility is $\sigma e^{-\kappa\tau}$ plus the long-run factor's.
\iqlookfor{mean reversion and a two-factor curve.}
\end{solution}

\begin{solution}{iq:m3:forward-curves-storage-and-convenience-yield:4}
\omterm{def:m3:forward-curves-storage-and-convenience-yield:floating}{Floating storage}: buy the near, charter a tanker, sell the far. Risks: charter rates rising with demand,
financing, operational losses, and the spread narrowing before you can lock it.
\iqlookfor{the carry trade and its costs.}
\end{solution}

\begin{solution}{iq:m3:forward-curves-storage-and-convenience-yield:5}
Event study: spread returns on the window's days against other days, controlling for curve level and
seasonality, over many months; test whether the effect scales with index size; placebo windows.
\iqlookfor{an event study with controls and placebos.}
\end{solution}

\begin{solution}{iq:m3:forward-curves-storage-and-convenience-yield:6}
Roll date (expiry, volume switch, fixed days), adjustment (none, difference, ratio), which contract to
hold, and the handling of negative prices. Traps: ratio adjustment across negative prices, back-adjusted
prices that no one traded, and signals computed on adjusted levels.
\iqlookfor{roll rules, adjustment choices and their artefacts.}
\end{solution}
